\uuid{M4XN}
\exo7id{2684}
\titre{exo7 2684}
\auteur{matexo1}
\organisation{exo7}
\datecreate{2002-02-01}
\isIndication{false}
\isCorrection{false}
\chapitre{Fonction de plusieurs variables}
\sousChapitre{Autre}
\cpge{mpsi-mp2i.fonctions-deux-variables.d, pcsi.fonctions-deux-variables.d, ptsi.fonctions-deux-variables.d, bcpst1.fonctions-deux-variables.a, mp-mpi.calcul-differentiel-optimisation.f}
\module{Analyse}
\niveau{L2}
\difficulte{}

\contenu{
\question{
Soit la fonction $f$ d\'efinie sur $\R_+^n$
par $f(x_1,\dots,x_n) = \sqrt[n]{x_1\ldots x_n}$. Chercher ses
extrema quand $x_1+\cdots+x_n = n$. En d\'eduire que la moyenne
g\'eom\'etrique de $n$ nombres positifs est plus petite que leur moyenne
arithm\'etique.
}
}
